% Generated draft: requires mathematical and source review.
\documentclass[11pt]{article}
\usepackage[margin=25mm]{geometry}
\usepackage{fontspec}
\setmainfont{Latin Modern Roman}
\usepackage{amsmath,amssymb,mathtools,unicode-math}
\setmathfont{Latin Modern Math}
\usepackage{enumitem,needspace,xurl,hyperref}
\setlength{\parindent}{0pt}
\setlength{\parskip}{5pt}
\setlength{\emergencystretch}{3em}
\newcommand{\N}{\mathbb{N}}
\newcommand{\Z}{\mathbb{Z}}
\newcommand{\Q}{\mathbb{Q}}
\newcommand{\R}{\mathbb{R}}
\newcommand{\C}{\mathbb{C}}
\newcommand{\F}{\mathbb{F}}
\newcommand{\A}{\mathbb{A}}
\newcommand{\PP}{\mathbb{P}}
\newcommand{\E}{\mathbb{E}}
\newcommand{\eps}{\varepsilon}
\newcommand{\dd}{\,\mathrm{d}}
\begin{document}
\section*{UnsolvedMath ID: 20002709}
\textbf{Source problem code:} catalogue entry
\subsection{Definitions and mathematical statement}
Fix integers $n,k\ge1$ and the half-turn on a regular polygon. The source requests elementary explanations for two type-$B$ proportions.

First, on a $(2kn+2)$-gon, consider centrally symmetric noncrossing partial $(k+2)$-angulations. An internal chord is $k$-admissible if it can occur in a full $(k+2)$-angulation; equivalently its two resulting polygons have vertex counts congruent to $2$ modulo $k$. Count a diameter as a one-element half-turn orbit and other chords in two-element orbits. Among partial angulations with exactly $i$ chord orbits, $0\le i\le n$, S. Fomin's asserted diameter-containing proportion is $i/n$. At most one diameter can occur in a noncrossing collection.

Second, on a $2kn$-gon, consider centrally symmetric noncrossing partitions whose nonzero blocks have cardinalities divisible by $k$. A zero block is invariant under the half-turn, equivalently contains an antipodal pair, and is unique when present. Its size is automatically a multiple of $2k$. Other blocks occur in antipodal pairs. The printed D. Armstrong claim says that among partitions with $b$ nonzero block-pairs the zero-block proportion is $b/n$. The supplied assessment says this is the complement of the correct proportion: it is
\[
 \frac{n-b}{n},\qquad 0\le b\le n,
\]
while $b/n$ is the zero-free proportion. Already $b=0$ gives the unique all-vertices zero block and refutes the printed claim, which would predict probability zero. The indexing must not be repaired silently by replacing the stated block count with its complement.

The aid reports the exact counts. Let $F_i$ count the partial angulations with $i$ chord orbits and $D_i$ those containing a diameter; let $H_b$ count the partitions with $b$ nonzero pairs and $Z_b$ those containing a zero block. Then
\[
 F_i=\binom ni\binom{kn+i}i,\qquad D_i=\frac inF_i,
\]
\[
 H_b=\binom nb\binom{kn}b,\qquad Z_b=\frac{n-b}{n}H_b.
\]
For $F(x)=\sum_{i=0}^nF_ix^i$, $D(x)=\sum_iD_ix^i$, $H(x)=\sum_bH_bx^b$ and $Z(x)=\sum_bZ_bx^b$, it further gives the positive pointed $f/h$ identity
\[
 F(x)-D(x)=(1+x)^{n-1}Z\left(\frac{x}{1+x}\right),
\]
or coefficientwise
\[
 F_i-D_i=\sum_{b=0}^{\min(i,n-1)}
               \binom{n-1-b}{i-b}Z_b.
\]
Thus the non-diameter face enumerator connects to the zero-block enumerator after the complement correction. The source's original remarks identify the first counts with $f$-numbers of $\Delta^{(k)}_{B_n}$ and the second with its $h$-numbers, describe the similarity as striking, and say that no connection with $NC^{(k)}_{B_n}$ was then known. That is historical problem context, not a present-day literature conclusion. The aid's exact-count and transform claims do not by themselves supply a bijective elementary explanation beyond their stated summation proof, and the broad request for elementary proofs is retained separately.
\subsection{Short English statement}
The diameter proportion is $i/n$, but the printed zero-block proportion is complemented: it is $(n-b)/n$ for $b$ nonzero block-pairs. The aid records exact counts and a positive $f/h$ transform linking non-diameter faces to zero-block partitions.
\subsection{Sources}
\begin{itemize}
\item American Institute of Mathematics workshop \emph{Braid groups, clusters and free probability}, Problem 2.3, parts attributed to S. Fomin and D. Armstrong, \url{https://aimath.org/WWN/braidgroups/braidgroups.pdf}.
\item Source reference labels [21, 51] for admissible chords and angulations, [1, 42] for divisible noncrossing partitions, and [21] in the $f/h$ remark. Full bibliographic entries are not supplied; the remark's trailing ``See'' is incomplete.
\item Local UnsolvedMath registry, permanent ID 20002709, source code \texttt{catalogue entry}, canonical location \texttt{aim-probability-notes.json}, \texttt{notes[150]}, record dated 14 August 2026. Statement status is reconstructed unverified, clean statement null, and original literature field empty.
\item Supplied research attempt 1: counterexample, problem status at run refuted, rigor labelled proved, novelty confidence low. Its complement correction, fixed-type summation and pointed transform are machine-generated research aid, not peer reviewed; independently verify before citation.
\end{itemize}
\end{document}
